Can a frozen code language model use execution feedback through a compact learned belief, without updating the actor's weights? We propose Particle-Belief Program Filtering (PBPF), which separates feedback interpretation from repair generation. A small belief module maintains candidate-specific uncertainty over latent failure mechanisms and predicts future execution outcomes. Its particles are updated sequentially on unchanged code, with importance-corrected transitions reserved for newly patched versions. A learned projector converts one sampled particle into soft-prefix tokens that remain fixed during an entire repair. The actor is frozen, but both the belief and its conditioning interface require training; the method therefore does not claim to be training-free. We define a prospective study with a primary seven-billion-parameter code actor, two model replications, source-disjoint calibration, and evaluator-hidden tests. The protocol first requires predictive value beyond matched deterministic memories before assessing repair utility using shared initial candidates, four repair calls, and complete token and execution accounting. It also distinguishes generation cost from belief computation and treats full-ensemble decoding as additional actor work. No confirmatory measurements are available in this package, so accuracy, resource costs, and replication findings remain pending. The study asks whether a compact explicit posterior can support useful feedback-grounded control, while making its training requirements, failure conditions, and information boundaries auditable.
